\documentclass[10pt,a4paper]{amsart}
\usepackage[margin=19mm]{geometry}
\usepackage{amsmath,amssymb,mathtools}
\usepackage[scr=rsfs,cal=euler]{mathalfa}
\usepackage{xfrac}
\usepackage[unicode,hidelinks,bookmarks=false]{hyperref}
\newcommand{\ie}{i.e.\ }
\newcommand{\cf}{cf.\ }
\newcommand{\resp}{respectively\ }
\newcommand{\Subsection}{\subsection}
\providecommand{\nobreakdash}{\nobreak\hbox{-}}
\setlength{\emergencystretch}{2em}
\multlinegap0pt
\usepackage{bm}
\usepackage{bbm}
\usepackage{dsfont}
\usepackage{xspace}
\usepackage{cancel}
\usepackage{mathtools}
\usepackage{amsthm}
\makeatletter
\@ifundefined{lemma}{\newtheorem{lemma}{Lemma}}{}
\makeatother
\newcommand{\B}{\mathscr{B}}
\newcommand{\C}{{\mathcal C}}
\newcommand{\Domain}{\mathcal{D}}
\newcommand{\Nc}{{\mathcal N}}
\newcommand{\R}{{\mathbb R}}
\renewcommand{\P}{\mathbb{P}}
\title[Finding the convex envelope of a boundary datum using random]{Finding the convex envelope of a boundary datum using random geometric graphs}
\author{Aurelia Deshayes, Nicol\'as Frevenza, Alfredo Miranda, Julio D. Rossi}
\date{Source: arXiv:2602.07696v1 (CC BY 4.0)}
\begin{document}
\maketitle

\noindent\emph{Source and licence.} Finding the convex envelope of a boundary datum using random geometric graphs. Version arXiv:2602.07696v1 (CC BY 4.0), distributed under the Creative Commons Attribution 4.0 International licence, as recorded in the arXiv licence metadata for this exact version and shown on the abstract page https://arxiv.org/abs/2602.07696v1. Full rights notice: licenses/arxiv-2602.07696-CC-BY-4.0.txt.\par\smallskip

\noindent\textit{Editorial scope.} Counted unit: the Lemma whose source label is \texttt{lema-cp} in the article (source0.tex lines 736--742, source label \texttt{lema-cp}), delivered together with the article's own complete original proof of it (source0.tex lines 744--855, one \texttt{proof} environment of 112 lines). No mathematics is written, repaired or re-derived: statement and proof are the article's own text line for line. Required context retained uncounted: DPP.intro (lines 426--442); condicion.lemma1 (lines 563--567); lemma-todas-las-dir (lines 655--661); lema-juego-termina (lines 686--691). Every definition, display and label of those lines is retained unchanged. Omitted from this package and not redistributed: 1--425, 443--562, 568--654, 662--685, 692--735, 743--743, 856--1625. Nothing inside a retained excerpt is omitted or altered: every retained body is delivered exactly as the article's lines are written, so no omission receipt of any kind accompanies this package. Citations: the retained excerpts carry no citation command at all, so no bibliography entry of the article is transcribed and no reference is redistributed. Reference adaptation: none was needed and none was made; every reference inside the delivered unit resolves inside it, so no unresolved reference or citation is produced. Printed slips kept literal: no printed word, symbol, hypothesis or number of the counted statement or of its proof was altered. Layout and class: the article's own class is \texttt{amsart} with packages including latexsym, amssymb, amsmath, dsfont, fancybox, xcolor, inputenc, amsfonts, systeme, mathtools, amsthm, graphicx, mathrsfs, enumitem; the wrapper is the lane's pinned \texttt{amsart} editorial wrapper at 10pt on A4, which restates the theorem-like environments the retained text needs and adds the article's own macro definitions that the retained text needs (\texttt{\textbackslash B}, \texttt{\textbackslash C}, \texttt{\textbackslash Domain}, \texttt{\textbackslash Nc}, \texttt{\textbackslash P}, \texttt{\textbackslash R}), with extra wrapper packages (bm, bbm, dsfont, xspace, cancel, mathtools). The delivered numbering is the unit's own. Assets: no retained span contains a figure environment, a graphics-inclusion command, a PDF-page inclusion, a table or a TikZ picture; no image file is copied into this package, the package declares no asset dependency and no image file is redistributed with it. Licence: version arXiv:2602.07696v1 of this article is distributed under the Creative Commons Attribution 4.0 International licence, as recorded in the arXiv licence metadata for this exact version and shown on the abstract page https://arxiv.org/abs/2602.07696v1. A full rights notice is delivered at licenses/arxiv-2602.07696-CC-BY-4.0.txt. Characters: every retained excerpt is pure ASCII, so the delivered engine, XeLaTeX with its default Latin Modern text font, reports no missing character.
\begin{equation}
\label{DPP.intro}
\begin{cases}
\displaystyle u(x)=\min\limits_{y\in \Nc_x^{\delta_n}} 
\left(
\frac{1}{2} 
u(y) 
+
\frac{1}{2}
u(y_x) 
\right)
& x\in \Domain_n \\
u(x)=f(x)
& x\in \B_n
\end{cases} 
.
\end{equation}

\begin{equation}
\label{condicion.lemma1}
\frac{c\sqrt{n} r_n^d \delta_n \alpha_n^{d-1}}{\log n} 
= \frac{\sqrt{n} \gamma_n}{\log n}\geq 2.
\end{equation}

\begin{lemma}
\label{lemma-todas-las-dir}
Assume condition~\eqref{condicion.lemma1}.
Then $\P-$almost surely there exists $n_0$, depending on the realization, such that for all $n \ge n_0$ there are points in each set of $\mathcal{A}_n$. 
Hence, $\Nc_x^{\delta_n}$ contains points in ``all directions" for every $x \in \C_n$, for $n\ge n_0$. 
In particular, $\Nc_x^{\delta_n} \neq \emptyset$ for all $x \in \C_n$.
\end{lemma}

\begin{lemma}
\label{lema-juego-termina}
Let $\tau$ be the hitting time of $\B_n$.
Suppose that $n\geq n_0,$ where $n_0$ is given by Lemma \ref{lemma-todas-las-dir}.
Then $$P_{S,n}^{x_0}(\tau < \infty)=1,$$ for any initial position $x_0\in\Domain_n$ and strategy $S$.
\end{lemma}

\begin{lemma}[Comparison Principle]
\label{lema-cp} 
Assume that $n\geq n_0$, where $n_0$ is given by Lemma \ref{lemma-todas-las-dir}.
Let $f, g \colon D^c \cap [0,1] \to \R$ be bounded functions with $f \leq g$.
Let $u$ be a subsolution of the DPP \eqref{DPP.intro} with boundary data $f$ on $\B_n$, and let $v$ be a supersolution of the DPP that coincides with $g$ on $\B_n$.
Then $u \leq v$.
\end{lemma}

\begin{proof}
We argue by contradiction. 
Assume that there exists some vertex $x_0\in\C_n$ such that $u(x_0)>v(x_0)$. 
Then,
\begin{equation}
M
=
\max_{y\in \C_n}
(u-v)(y)>0,
\end{equation}
and it is attained at some $x\in \Domain_n$.
Using that $u$ and $v$ are sub and super-solutions of the DPP, we have that,
\begin{align*}
M &
= u(x)-v(x) \\
&\displaystyle  \leq 
\min_{y\in \Nc_x^{\delta_n}} 
\left(
\frac{1}{2} 
u(y) 
+
\frac{1}{2}
u(y_x) 
\right)
-
\min_{y\in \Nc_x^{\delta_n}} 
\left(
\frac{1}{2} 
v(y) 
+
\frac{1}{2}
v(y_x) 
\right)
\\
& \displaystyle 
\leq 
\frac{1}{2} 
u(z) 
+
\frac{1}{2}
u(z_x) 
-
\frac{1}{2} 
v(z) 
-
\frac{1}{2}
v(z_x)
\\
&\displaystyle  =
\frac{1}{2} 
\left(
u(z) 
-
v(z) 
\right)
+
\frac{1}{2}
\left(
u(y_x) 
-
v(y_x)
\right)
\\
& \displaystyle  \leq M
,
\end{align*}
where $z\in \Nc_x^{\delta_n}$ is such that
\[
\min_{y\in \Nc_x^{\delta_n}} 
\left(
\frac{1}{2} 
v(y) 
+
\frac{1}{2}
v(y_x) 
\right)
=
\frac{1}{2} 
v(z) 
+
\frac{1}{2}
v(z_x)
.
\]
It follows that 
\[
u(z) 
-
v(z) 
= M
\qquad 
\text{and}
\qquad
u(z_x) 
-
v(z_x)
=
M
.
\]
Using the same argument given in Lemma \ref{lema-juego-termina}, where with probability 1/2 the distance to a fixed point increases at each step by at least a fixed amount, we have that eventually reachs the set $\B_n$,
Hence, for some $y\in \B_n$ we obtain
\[
u(y)-v(y) 
= 
M
>
0,
\]
which is a contradiction, since we have $$u(y)=f(y)\leq g(y)= v(y)$$
for all $y\in\B_n$.
\end{proof}

\end{document}
